% =============================================================================
% CHAPTER 22: POLICY APPLICATIONS AND SECTOR-SPECIFIC PORTFOLIO DESIGN
% =============================================================================
% HEADER BLOCK
% -----------------------------------------------------------------------------
% Document:      Chapter 22
% Category:      Application (Type C)
% Version:       1.0 (January 2026)
% Dependencies:  Chapters 16-21, Appendices HHH, KKK, JJJ, BBB, IIE
% -----------------------------------------------------------------------------
%
% Purpose: Demonstrate how Chapters 16-21 (Status Quo Diagnosis, Single
% Interventions, Phase-Affinity, Segment Multipliers, Static Portfolio Design,
% Dynamic Portfolio Evolution) apply to major policy domains.
%
% This chapter bridges from theory to practice by showing real policy portfolios
% across Health, Pensions, Environmental, and Labor Market domains, and analyzing
% their coherence and interaction effects.
%
% Primary Appendix: HHH (METHOD-TOOLKIT) - Complementarity matrix, coherence
% Secondary Appendices: KKK (METHOD-DYNAMIC), JJJ (METHOD-PSG), BBB, IIE
%
% Prerequisites: Chapters 16-21 (Status Quo → Single Interventions →
% Phase-Affinity → Segment-Multipliers → Static Design → Dynamic Evolution)
% Continues from: Chapter 21 (Dynamic Portfolio Evolution)
% Leads to: Chapter 23 (Multi-Level Implementation)
%
% Chapter Type: C (Application)
%
% =============================================================================

\chapter{Policy Applications and Sector-Specific Portfolio Design}
\label{ch:policy-applications}

%%%%%%%%%%%%%%%%%%%%%%%%%%%%%%%%%%%%%%%%%%%%%%%%%%%%%%%%%%%%%%%%%%%%%%%%%%%%%%%
% CHAPTER SCOPE BOX
%%%%%%%%%%%%%%%%%%%%%%%%%%%%%%%%%%%%%%%%%%%%%%%%%%%%%%%%%%%%%%%%%%%%%%%%%%%%%%%
\begin{tcolorbox}[
    colback=purple!5!white,
    colframe=purple!75!black,
    title={\textbf{Chapter Scope: Ziel / In-Scope / Out-of-Scope}},
    fonttitle=\bfseries
]

\textbf{Ziel (Objective):}
Apply the EBF framework to real-world implementation with practical examples and policy implications.

\vspace{0.3cm}
\textbf{In-Scope:}
\begin{itemize}[nosep]
    \item Practical applications of real-world implementation
    \item Multiple worked examples (≥2)
    \item Policy implications and recommendations
    \item Implementation guidance
\end{itemize}

\vspace{0.3cm}
\textbf{Out-of-Scope (delegiert an Appendices):}
\begin{itemize}[nosep]
    \item Theoretical foundations $\rightarrow$ Foundation chapters
    \item Formal proofs $\rightarrow$ FORMAL appendices
    \item Parameter estimation methods $\rightarrow$ METHOD appendices
\end{itemize}

\end{tcolorbox}



%%%%%%%%%%%%%%%%%%%%%%%%%%%%%%%%%%%%%%%%%%%%%%%%%%%%%%%%%%%%%%%%%%%%%%%%%%%%%%%
% CHAPTER OVERVIEW
%%%%%%%%%%%%%%%%%%%%%%%%%%%%%%%%%%%%%%%%%%%%%%%%%%%%%%%%%%%%%%%%%%%%%%%%%%%%%%%
\begin{tcolorbox}[
    colback=green!5!white,
    colframe=green!75!black,
    title={\textbf{Chapter Overview}},
    fonttitle=\bfseries
]
This chapter demonstrates how the Ch16-21 portfolio framework applies to real policy domains:

\begin{enumerate}[nosep]
    \item \textbf{Section \ref{sec:ch22-intro}}: Introduction---from individual behavior to policy design
    \item \textbf{Section \ref{sec:health-portfolio}}: Health policy portfolio design (diabetes prevention)
    \item \textbf{Section \ref{sec:pension-portfolio}}: Pension policy portfolio design (reform acceptance)
    \item \textbf{Section \ref{sec:environmental-portfolio}}: Environmental policy portfolio (energy reduction)
    \item \textbf{Section \ref{sec:labor-portfolio}}: Labor market policy (organizational engagement)
    \item \textbf{Section \ref{sec:cross-domain}}: Cross-domain coherence and spillovers
    \item \textbf{Section \ref{sec:implementation}}: Practical implementation framework
    \item \textbf{Section \ref{sec:ch22-summary}}: Summary and conclusions
\end{enumerate}
\end{tcolorbox}



% =============================================================================
% QUICK REFERENCE BOX
% =============================================================================

\begin{tcolorbox}[colback=blue!5!white, colframe=blue!75!black, title=Key Concepts in This Chapter]
\small
\begin{tabular}{@{}ll@{}}
\textbf{Policy Portfolio} & Set of coordinated interventions in a domain \\
\textbf{Domain} & Health, Pensions, Environment, Labor Markets \\
\textbf{S_0 (Sector Status)} & Current state of population in domain \\
\textbf{Policy Coherence} & Do interventions reinforce or interfere? \\
\textbf{Spillovers} & How do portfolios in one domain affect others? \\
\textbf{Implementation Sequence} & Which policy changes need to happen first? \\
\textbf{Measurement} & How to track policy portfolio effectiveness? \\
\end{tabular}
\end{tcolorbox}

% =============================================================================
% APPENDIX REFERENCES
% =============================================================================

\begin{tcolorbox}[colback=yellow!10!white, colframe=orange!75!black, title=Appendix References]

\textbf{Essential:}
\begin{itemize}[nosep]
\item Appendix TKT (METHOD-TOOLKIT) --- Complementarity matrix, coherence criteria
\item Appendix UNMAPPED_KKK (METHOD-DYNAMIC) --- Formal models for dynamic adjustment
\item Appendix UNMAPPED_LLL (METHOD-POLICY-DOMAINS) --- Domain-specific portfolio templates
\item Appendix UNMAPPED_MMM (METHOD-POLICY-COHERENCE) --- Cross-domain spillover analysis
\end{itemize}

\textbf{Supporting:}
\begin{itemize}[nosep]
\item Appendix UNMAPPED_JJJ (METHOD-PSG) --- Parameter sourcing for policy contexts
\item Appendix UNMAPPED_BBB (CORE-WHERE) --- Parameter repository
\item Appendix IIE (PREDICT-IMPACT) --- Impact forecasting for policy interventions
\item \textbf{Appendix UNMAPPED_RB (LIT-BENABOU)} --- Narrative-aware policy design ($I_{\text{AWARE}}$ + framing)
\end{itemize}

\end{tcolorbox}

% =============================================================================
% CROSS-REFERENCE: EMERGE ALGORITHM (Chapter 11, EIT-EMP)
% =============================================================================
\begin{tcolorbox}[colback=blue!5!white, colframe=blue!75!black,
    title=\textbf{Cross-Reference: EMERGE Algorithm \& Policy-Level K* (Chapter 11, 20)}]

Policy-Portfolios operieren auf \textbf{Scope Level L0-L1} (Systemic/Strategic) des EMERGE-Algorithmus (Kapitel 11, EIT-EMP-3, \textbf{EXC-3 Hybrid}):

\vspace{0.2cm}
\textbf{Policy-Level K* (EIT-EMP-2, EXC-3):}
\begin{equation*}
K^*_{\text{policy}} = K_{\text{base}} \cdot f_B + \delta_C + \delta_S + \delta_\gamma + \underbrace{\delta_L}_{\text{L0: -40\%, L1: -20\%}}
\end{equation*}
\textit{Nur $f_B$ multiplikativ; $\delta_L$ ist additiver Scope-Adjustment.}

\textbf{Scope-Level Adjustment für Policy:}
\begin{itemize}[nosep]
\item \textbf{L0 (Systemic):} $f_L = 0.6$ → $K^* \approx 3-4$ (wenige, aber tiefgreifende Interventionen)
\item \textbf{L1 (Strategic):} $f_L = 0.8$ → $K^* \approx 4-6$ (policy bundles pro Domain)
\end{itemize}

\vspace{0.2cm}
\textbf{Policy-Constraints (Ch20):}
\begin{itemize}[nosep]
\item C6 (Segment): $\sigma_s \geq 0$ für alle Major-Segmente (kein Backlash bei 25\%+ der Population)
\item C7 (K-Optimal): $|P_{\text{domain}}| \leq K^*_{\text{policy}}$ (Fokus statt Streuung)
\item C2 (Crowding): $\gamma_{ij} \geq -0.1$ zwischen Policy-Domains (Cross-Domain Coherence)
\end{itemize}

\vspace{0.2cm}
\textit{Siehe: Chapter 11 §11.20 (EMERGE), Chapter 20 §20.7 (Constraints), Appendix UNMAPPED_IE (EMERGE-SPEC)}
\end{tcolorbox}

% =============================================================================
% INTUITION BOX
% =============================================================================

\begin{tcolorbox}[colback=yellow!5!white, colframe=yellow!75!black, title=Intuition: Anna's Government Policy Challenge]

\textbf{Anna}, Chief of Health Policy in a European region, faces a puzzle. Individual interventions work:
\begin{itemize}[nosep]
\item Information campaigns increase health awareness (+8-12\%)
\item Incentive programs increase medication adherence (+15-20\%)
\item Social programs improve mental health (+10-15\%)
\end{itemize}

But when she launches \emph{all three together}, the combined effect is only +18\%, not the expected +30-35\%. Why?

\textbf{The Hidden Problem:} The interventions work on different segments with different effectiveness, create contradictory messaging, and compete for limited implementation capacity. Worse: the public message ``we're fixing health'' without clear priority creates cynicism in exactly the Present-Biased segment most likely to drop out.

\textbf{What Anna needs:} A framework showing which interventions synergize (+γ), which interfere (-γ), and how to sequence them over time so that early wins build toward later changes. Chapters 16-21 provided the theory. This chapter shows how to apply it.

\end{tcolorbox}

% =============================================================================
% CENTRAL QUESTION BOX
% =============================================================================

\begin{tcolorbox}[colback=red!5!white, colframe=red!75!black, title=Central Question of This Chapter]

\textbf{How do we apply the Chapters 16-21 portfolio framework to real policy domains?}

\vspace{0.3cm}

More specifically:
\begin{itemize}
\item How do we diagnose policy domain Status Quo (using Ch16 framework)?
\item Which intervention vectors $\vec{I} \in [0,1]^9$ make sense in each domain (using Ch17)?
\item How do phase-type and segment-type multipliers change domain-to-domain (using Ch18-19)?
\item What static portfolio emerges when we optimize for the domain context (using Ch20)?
\item How does the portfolio evolve as the domain conditions change (using Ch21)?
\item When do policy portfolios in different domains interfere with each other?
\end{itemize}

This chapter demonstrates all six questions across four major policy domains.

\end{tcolorbox}

---

% =============================================================================
% MAIN CONTENT
% =============================================================================

\section{Introduction: From Individual Behavior to Policy Design}
\label{sec:ch22-intro}

Chapters 16-21 developed a complete framework for portfolio design starting from individual behavior change. Each chapter focused on a single person or organization making decisions:

\begin{itemize}
\item \textbf{Ch16:} Diagnose individual Status Quo (the 9C vector)
\item \textbf{Ch17:} Design emergent intervention vectors ($\vec{I} \in [0,1]^9$)
\item \textbf{Ch18-19:} Modulate effectiveness by phase and segment
\item \textbf{Ch20:} Optimize static portfolio given status quo
\item \textbf{Ch21:} Redesign portfolio as conditions evolve
\end{itemize}

Policy domains operate at a different scale---we care about population-level effects, not individual outcomes. Yet the same logic applies:

\begin{quote}
\textbf{Definition:} A \emph{policy portfolio} is a set of coordinated interventions designed to change population-level behavior in a domain, subject to coherence constraints and interaction effects.
\end{quote}

The key insight: Portfolio design at the policy level follows the same principles as individual-level design, with three modifications:

\begin{enumerate}
\item \textbf{Scale-up of segment heterogeneity:} Populations include diverse segments (Present-Biased, Social-Oriented, Autonomy-Seeking). Effective policies target each segment differently.

\item \textbf{Implementation dynamics:} Policy changes are constrained by administrative capacity, political feasibility, and path-dependency. Sequencing matters.

\item \textbf{Cross-domain spillovers:} Policies in one domain (health) affect outcomes in others (labor markets, retirement security). Coherence requires managing these spillovers.
\end{enumerate}

This chapter walks through four policy domains---Health, Pensions, Environment, and Labor Markets---showing how Ch16-21 framework applies to each, and where interactions emerge.

% =============================================================================
% SECTION 1: HEALTH POLICY PORTFOLIO
% =============================================================================

\section{Domain 1: Health Policy Portfolio Design}
\label{sec:health-portfolio}

\subsection{Status Quo Diagnosis (Using Ch16 Framework)}

Consider a regional health authority designing interventions for Type-2 diabetes prevention in a population of 500,000 adults.

Using Chapter 16's Status Quo diagnostic:

\begin{center}
\begin{tabular}{@{}lcc@{}}
\toprule
\textbf{9C Dimension} & \textbf{Current State} & \textbf{Target State} \\
\midrule
\textbf{L (Who)} & General adults, mixed ages & Focus on age 40-65 \\
\textbf{d (What)} & Multiple: Health, Social, Financial & Primary: Health (u\_H) \\
\textbf{γ (How)} & Many comorbidities, low social support & Reduced interaction effects \\
\textbf{Ψ (Context)} & Low health literacy, high obesity rate & Supportive environment \\
\textbf{Θ (Parameters)} & Unknown effectiveness of interventions & Estimated via pilots \\
\textbf{A (Awareness)} & 18\% aware of diabetes risk & 60\% target \\
\textbf{W (Readiness)} & 8\% ready to change behavior & 35\% target \\
\textbf{φ (Phase)} & Mostly PreAware (65\%), some Triggered (25\%) & Shift to Action phase \\
\textbf{Scope} & Individual level, no community support & Multi-level (individual + org + community) \\
\midrule
\end{tabular}
\end{center}

Key insight: The population is predominantly in \textbf{PreAware phase} (φ=0.2). This constrains which interventions will work (see Phase-Affinity, Ch18).

\subsection{Single Intervention Assessment (Using Ch17 Framework)}

From the 10C intervention framework (Chapter 17), what emergent intervention vectors are available?\footnote{Interventions \textit{emerge} from the continuous 10C space $[0,1]^9$---they are not selected from a catalogue (see Ch.~17, Axiom EIT-3). The labels below are shorthand for communication.}

\begin{center}
\begin{tabular}{@{}lllcc@{}}
\toprule
\textbf{10C Emphasis} & \textbf{Intervention} & \textbf{Mechanism} & \textbf{Feasibility} & \textbf{Est. Effect} \\
\midrule
AWARE & Health information campaign & Increase A & High & 8-12\% \\
READY & Screening + feedback app & Self-correction & High & 12-15\% \\
WHEN & Workplace-based choice arch. & Default options & Medium & 10-14\% \\
STAGE & Time-limited intervention push & Urgency & Low & 5-8\% \\
WHO & Identity-based messaging & Self-concept & Medium & 6-10\% \\
WHAT$_S$ & Community peer support groups & Social norms & Medium & 15-20\% \\
WHAT$_F$ & Subsidy for gym memberships & Financial & Low & 4-8\% \\
HOW & Pre-commitment contracts & Binding intent & Medium & 18-25\% \\
\midrule
\end{tabular}
\end{center}

All types are available, but effectiveness varies by phase and segment (next sections).

\subsection{Phase-Type Affinity (Using Ch18 Framework)}

Recall from Chapter 18: Phase-type affinity α measures how well intervention type matches BCJ phase.

For \textbf{PreAware phase} (φ=0.2), which interventions have high affinity?

\begin{center}
\begin{tabular}{@{}llcc@{}}
\toprule
\textbf{10C Emphasis} & \textbf{Intervention} & \textbf{α (Affinity)} & \textbf{Interpretation} \\
\midrule
AWARE & Information campaign & 0.85 & \textbf{Excellent} --- Information essential at this phase \\
READY & Feedback/screening & 0.70 & Good --- Can trigger transition to Triggered phase \\
WHEN & Choice architecture & 0.55 & Fair --- Weak when awareness is lacking \\
STAGE & Time-limited push & 0.40 & Poor --- Urgency ineffective for PreAware \\
WHO & Identity messaging & 0.60 & Fair --- Can work if identity is salient \\
WHAT$_S$ & Peer support & 0.50 & Fair --- Peer pressure weak without awareness \\
WHAT$_F$ & Financial incentives & 0.35 & Poor --- Money doesn't motivate without awareness \\
HOW & Commitment & 0.30 & Very Poor --- Can't commit to changing behavior you don't understand \\
\midrule
\end{tabular}
\end{center}

Result: For PreAware phase, \textbf{prioritize AWARE (Information) and READY (Feedback)}.

STAGE, WHAT$_F$, HOW have poor affinity at this phase and should be introduced later, after population moves to Triggered phase.

\subsection{Segment-Multiplier Effects (Using Ch19 Framework)}

Within the population, different behavioral segments respond differently. From Ch19, three segments:

\begin{center}
\begin{tabular}{@{}lrrrrr@{}}
\toprule
\textbf{Segment} & \textbf{Size} & \textbf{AWARE} & \textbf{READY} & \textbf{WHAT$_S$} & \textbf{HOW} \\
\midrule
\textbf{Present-Biased} & 40\% & 0.6× & 0.8× & 0.5× & 0.4× \\
\textbf{Social-Oriented} & 35\% & 0.9× & 0.9× & 1.4× & 0.9× \\
\textbf{Autonomy-Seeking} & 25\% & 1.1× & 1.0× & 0.6× & 1.2× \\
\midrule
\end{tabular}
\end{center}

Interpretation:
\begin{itemize}
\item \textbf{Information campaigns (AWARE):} Most effective for Autonomy-Seeking (+10\%), less for Present-Biased (6\%)
\item \textbf{Peer support (WHAT$_S$):} Highly effective for Social-Oriented (+20\%), weak for Autonomy-Seeking (6\%)
\item \textbf{Commitment (HOW):} Works well for Autonomy-Seeking (+22\%), much less for Present-Biased (4\%)
\end{itemize}

Key insight: \textbf{One-size-fits-all policy doesn't work}. We need segment-specific messaging and design.

\subsection{Static Portfolio Optimization (Using Ch20 Framework)}

Using Chapter 20's optimization logic with constraints:
- Total budget: 10 million EUR
- Administrative capacity: Can implement 4-5 interventions
- Must have high phase-affinity (α > 0.6)

Recommended Static Portfolio P₀ for PreAware → Triggered transition:

\begin{center}
\begin{tabular}{@{}llll@{}}
\toprule
\textbf{Intervention} & \textbf{Budget} & \textbf{Expected Effect} & \textbf{Role} \\
\midrule
AWARE: Health information campaign & 2.5M & +9.8\% & Primary driver (high α=0.85) \\
READY: Screening + mobile app & 3.0M & +13.2\% & Trigger transition (α=0.70) \\
WHAT$_S$: Peer support program (SO segment) & 2.0M & +8.5\% & Reinforce for Social-Oriented \\
WHO: Identity messaging (AS segment) & 1.5M & +6.1\% & Engage Autonomy-Seeking \\
\midrule
\end{tabular}
\end{center}

Combined Effect (using Coherence Formula from Ch20, Eq. 906):
$$E(P_0) = \sum_i E_i \cdot \alpha_i \cdot \sigma_s + \sum_{i<j} \gamma_{ij} \sqrt{E_i \cdot E_j}$$

With complementarity estimates:
- γ(AWARE, READY) = +0.3 (information + feedback reinforce)
- γ(AWARE, WHAT$_S$) = +0.2 (information + peer support reinforce)
- γ(WHO, WHAT$_S$) = -0.1 (identity-based + peer messaging can conflict slightly)

\textbf{Net Effect: +40-45\% population shift toward Triggered phase within 12 months.}

\subsection{Dynamic Portfolio Adjustment (Using Ch21 Framework)}

As population moves to Triggered phase (φ increases), which interventions should we adjust?

After 6 months (Ch21 dynamics):

\begin{center}
\begin{tabular}{@{}llccc@{}}
\toprule
\textbf{Phase} & \textbf{Month} & \textbf{Redesign Trigger} & \textbf{Add} & \textbf{Remove} \\
\midrule
PreAware & t=0 & --- & P₀ portfolio & --- \\
Triggered & t=6 & φ increased to 0.55 & STAGE (urgency), WHAT$_F$ (incentives) & Reduce AWARE (info redundant) \\
Action & t=12 & φ increased to 0.75 & HOW (commitment), WHEN (choice arch) & Reduce READY (less feedback needed) \\
Maintenance & t=18+ & φ reaches 0.90 & WHAT$_S$ boost (ongoing support) & --- \\
\midrule
\end{tabular}
\end{center}

The logic: Different phases require different intervention mixes. As population evolves, redesign the portfolio to maintain effectiveness.

% =============================================================================
% SECTION 2: PENSION/RETIREMENT POLICY PORTFOLIO
% =============================================================================

\section{Domain 2: Pension Policy Portfolio Design}
\label{sec:pension-portfolio}

\subsection{Status Quo Diagnosis for Pension Reform}

Consider a national government needing to reform retirement security due to demographic shift and declining birth rates. Using Ch16 diagnostic:

\begin{center}
\begin{tabular}{@{}lcc@{}}
\toprule
\textbf{Dimension} & \textbf{Current} & \textbf{Policy Target} \\
\midrule
\textbf{L (Who)} & Mostly retirees 65+, some future retirees & All age cohorts aware \\
\textbf{d (What)} & Financial security, social status & All 6 dimensions (FEPSDE) \\
\textbf{Ψ (Context)} & Trust in government declining, media polarized & Rebuild institutional trust \\
\textbf{A (Awareness)} & 25\% understand reform necessity & 70\% understand \\
\textbf{W (Readiness)} & 12\% accept reform necessity & 50\% accept \\
\textbf{φ (Phase)} & PreAware (70\%), some Triggered (20\%) & Majority Triggered (60\%+) \\
\textbf{Scope} & Individual, limited community engagement & Multi-level: individual → family → community \\
\midrule
\end{tabular}
\end{center}

Key challenge: This is a \textbf{context-shifting change} (Ψ changes, trust erodes). The portfolio must address both behavioral change AND institutional credibility.

\subsection{Policy Portfolio for Pension Reform}

The reform portfolio uses Chapters 17-20 logic:

\begin{center}
\begin{tabular}{@{}lll@{}}
\toprule
\textbf{Intervention} & \textbf{Type (Ch17)} & \textbf{Mechanism} \\
\midrule
1. Transparent impact modeling & AWARE (Awareness) & Show individual consequences \\
2. Citizens' assembly on reform options & WHAT$_S$ (Social) & Build deliberative consensus \\
3. Staged benefit structure changes & STAGE (Temporal) & Implement over 10 years, not suddenly \\
4. Intergenerational framing & WHO (Identity) & Appeal to fairness for future generations \\
5. Flexible retirement pathway design & WHEN (Choice Arch) & Default + options for different preferences \\
6. Pre-commitment opt-in mechanisms & HOW (Commitment) & Early commitment for retirement savings \\
\midrule
\end{tabular}
\end{center}

Coherence Analysis (from HHH):

The six interventions must be coherent:
- Information (AWARE) must match actual policy (WHEN, STAGE)
- Citizen assembly (WHAT$_S$) must influence actual policy design (WHEN, WHO)
- Temporal staging (STAGE) must align with political feasibility

\textbf{Critical complementarity:}
\begin{itemize}
\item γ(AWARE, WHAT$_S$) = +0.4 --- Information + deliberation reinforce each other (high)
\item γ(WHEN, HOW) = +0.3 --- Choice architecture + opt-in reinforce (medium)
\item γ(WHO, STAGE) = +0.2 --- Identity framing + staged implementation reinforce (medium)
\item γ(WHO, WHAT$_S$) = +0.1 --- Identity framing + citizen input moderate positive
\end{itemize}

Expected static effect (using Ch20 formula): +52-58\% increase in acceptance of reform package, with support for specific implementation pathway.

\subsection{Dynamic Adjustment: Learning and Course Correction}

A 10-year reform requires dynamic adjustment (Ch21):

\begin{center}
\begin{tabular}{@{}lll@{}}
\toprule
\textbf{Year} & \textbf{Redesign Trigger} & \textbf{Action} \\
\midrule
Y0-Y3 & Information phase; monitor understanding & Continue AWARE, WHAT$_S$, add data transparency \\
Y3-Y6 & Implementation begins; watch trust trajectory & Increase WHAT$_S$ (ongoing deliberation), add WHAT$_F$ if needed \\
Y6-Y8 & Behavior changing; measure new retirement patterns & Monitor γ(STAGE,WHO) interaction (staging vs identity) \\
Y8-Y10 & Convergence phase; stabilize new equilibrium & Shift to WHAT$_S$ support (community cohesion) \\
\midrule
\end{tabular}
\end{center}

The framework allows adaptive management: If public trust drops during Y3-Y6, increase citizen engagement (WHAT$_S$) and transparency (AWARE). If behavioral change lags, strengthen commitment mechanisms (HOW).

% =============================================================================
% SECTION 3: ENVIRONMENTAL/CLIMATE POLICY PORTFOLIO
% =============================================================================

\section{Domain 3: Environmental Policy Portfolio Design}
\label{sec:environmental-portfolio}

\subsection{The Energia Worked Example Revisited}

Recall from Chapters 20-21 the Energia (energy consumption) worked example:

A city government seeks to reduce residential energy use 20\% within 5 years. Baseline: 2,500 kWh/year per household.

\subsection{Status Quo Analysis}

\begin{center}
\begin{tabular}{@{}lcc@{}}
\toprule
\textbf{Dimension} & \textbf{Current} & \textbf{Target} \\
\midrule
\textbf{Awareness} & 15\% understand carbon impact & 60\% understand \\
\textbf{Readiness} & 8\% ready to change & 40\% ready \\
\textbf{Phase} & 80\% PreAware, 15\% Triggered & 50\% Action, 20\% Maintenance \\
\textbf{Context} & Economic growth = energy growth norm & Energy efficiency as prosperity signal \\
\textbf{Segments} & 45\% Present-Biased, 30\% Social, 25\% Autonomy & All segments engaged \\
\midrule
\end{tabular}
\end{center}

\subsection{Static Portfolio Design}

Using Ch20 framework:

\begin{center}
\begin{tabular}{@{}llll@{}}
\toprule
\textbf{Strategy} & \textbf{Target Segment} & \textbf{Mechanism} & \textbf{Expected Effect} \\
\midrule
1. Public awareness campaign & All & AWARE (information) & 6\% \\
2. Home energy audits + feedback & Present-Biased & READY (feedback) & 8\% \\
3. Social comparison (energy reports) & Social-Oriented & WHAT$_S$ (social) & 10\% \\
4. Smart meter + financial incentive & Present-Biased & WHAT$_F$ (financial) & 7\% \\
5. Green building standards (choice arch) & Autonomy-Seeking & WHEN (choice) & 9\% \\
6. Community energy cooperatives & Social-Oriented & WHAT$_S$ (social) & 11\% \\
\midrule
\end{tabular}
\end{center}

Static Effectiveness: E(P₀) = 14\% average reduction (baseline: 2,500 kWh → 2,150 kWh).

But note: Different segments see different effects.

\subsection{Context Shock Scenario (Ch21 Dynamics)}

In Year 2, economic crisis hits. Energy prices drop; context Ψ changes:

\begin{quote}
\textbf{Scenario:} Recession, unemployment up 5pp, household budgets tight. Short-term thinking increases. Baseline effectiveness drops to 6\% because financial stress (Ψ) reduces readiness.
\end{quote}

Portfolio redesign required (using Ch21 framework):

\begin{center}
\begin{tabular}{@{}llcc@{}}
\toprule
\textbf{Intervention} & \textbf{Pre-Crisis Effect} & \textbf{Post-Crisis Effect} & \textbf{Redesign Action} \\
\midrule
Campaign + Feedback & 8\% & 3.2\% (40\% decline) & \textbf{Reduce scope, focus on basics} \\
Social comparison & 10\% & 4\% (60\% decline) & \textbf{Reframe as cost savings, not climate} \\
Financial incentive & 7\% & 9\% (increase!) & \textbf{Expand --- now more motivating} \\
Green standards & 9\% & 2\% (80\% decline!) & \textbf{Suspend --- triggers backlash in crisis} \\
\midrule
\end{tabular}
\end{center}

Redesigned Portfolio P₁ (crisis response):
- Keep awareness (AWARE): Low cost, still works
- Strengthen financial incentives (WHAT$_F$): Even more effective in recession
- Pause/reduce green standards (WHEN): High reactance risk in crisis
- New: Framing shift toward "cost savings" (WHO identity change) instead of "climate duty"

Result: Despite crisis, maintain 10-12\% effectiveness through portfolio adaptation.

Post-crisis recovery (Years 4-5): Shift back to broader portfolio as context stabilizes.

% =============================================================================
% SECTION 4: LABOR MARKET / ORGANIZATIONAL POLICY PORTFOLIO
% =============================================================================

\section{Domain 4: Labor Market Policy Portfolio Design}
\label{sec:labor-portfolio}

\subsection{Organizational Engagement During Disruption}

Recall from Chapters 20-21 the Engagement worked example:

A large manufacturing firm must maintain employee engagement during major restructuring (potential layoffs). Baseline morale: 15.6\%.

\subsection{Status Quo Diagnostic}

\begin{center}
\begin{tabular}{@{}lcc@{}}
\toprule
\textbf{Dimension} & \textbf{Baseline} & \textbf{Crisis (Post-Shock)} \\
\midrule
\textbf{Awareness} & 70\% understand company strategy & 40\% (information void) \\
\textbf{Readiness} & 55\% willing to adapt & 25\% (fear dominates) \\
\textbf{Phase} & 50\% Action, 30\% Maintenance & 70\% PreAware, 20\% Triggered \\
\textbf{Context} & Psychological safety high & Psychological safety collapsed \\
\textbf{Segments} & Mixed & Present-Biased dominates (anxiety) \\
\midrule
\end{tabular}
\end{center}

Key insight: Restructuring creates massive context shift. Interventions that worked pre-crisis now backfire.

\subsection{Baseline Portfolio (Pre-Crisis)}

\begin{center}
\begin{tabular}{@{}llcc@{}}
\toprule
\textbf{Intervention} & \textbf{Mechanism} & \textbf{Effect} & \textbf{Segment} \\
\midrule
1. Autonomy-focused feedback & READY: Empower decision-making & 0.30 × 0.20 = 6\% & Autonomy-Seeking \\
2. Peer recognition programs & WHAT$_S$: Social recognition & 0.40 × 0.15 = 6\% & Social-Oriented \\
3. Leadership communication & AWARE: Information + vision & 0.30 × 0.12 = 3.6\% & Present-Biased \\
\midrule
\textbf{Combined baseline} & & 15.6\% & --- \\
\midrule
\end{tabular}
\end{center}

\subsection{Crisis Shock and Collapse}

When restructuring is announced:

\begin{center}
\begin{tabular}{@{}llcc@{}}
\toprule
\textbf{Intervention} & \textbf{Pre-Crisis} & \textbf{Post-Crisis} & \textbf{Problem} \\
\midrule
1. Autonomy feedback & 0.30 × 0.20 = 6.0\% & 0.30 × 0.08 = 2.4\% & Autonomy seen as ``abandonment'' \\
2. Peer recognition & 0.40 × 0.15 = 6.0\% & 0.40 × 0.05 = 2.0\% & Unfair when layoffs coming \\
3. Leadership comms & 0.30 × 0.12 = 3.6\% & 0.30 × 0.04 = 1.2\% & Empty promises trigger cynicism \\
\midrule
\textbf{Combined post-crisis} & 15.6\% & 5.6\% & --- \\
\midrule
\end{tabular}
\end{center}

Collapse: 15.6% → 5.6% (64\% drop in 2 weeks).

\subsection{Redesigned Portfolio for Crisis Management}

Using Ch21 framework, completely different portfolio needed:

\begin{center}
\begin{tabular}{@{}llc@{}}
\toprule
\textbf{New Intervention} & \textbf{Mechanism} & \textbf{Expected Effect} \\
\midrule
1. Safety-first communication & STAGE: Urgency + AWARE: Honesty & ``Here's what's happening, here's the timeline'' \\
2. Job preservation focus & WHO: Identity shift & ``Protecting core teams; supporting transitions'' \\
3. Team cohesion support & WHAT$_S$: Social bonding & Peer support for those staying and leaving \\
4. Mental health services & WHEN: Access choice arch & Visible, accessible counseling \\
5. Commitment to transparency & HOW: Pre-commitment & ``We'll share information monthly, no surprises'' \\
\midrule
\end{tabular}
\end{center}

Redesigned effectiveness:

\begin{center}
\begin{tabular}{@{}llc@{}}
\toprule
\textbf{Strategy} & \textbf{Effect Size} & \textbf{Result} \\
\midrule
Safety communication & 0.40 × 0.12 = 4.8\% & Stabilizes anxiety \\
Job preservation messaging & 0.25 × 0.14 = 3.5\% & Shifts identity to ``survivor'' \\
Team support programs & 0.35 × 0.14 = 4.9\% & Restores social cohesion \\
Mental health access & 0.20 × 0.09 = 1.8\% & Enables individual coping \\
Transparency commitment & 0.20 × 0.10 = 2.0\% & Reduces fear from uncertainty \\
\midrule
\textbf{Combined redesigned} & & 12.1\% \\
\midrule
\end{tabular}
\end{center}

Recovery: 5.6% → 12.1% (115\% improvement through portfolio redesign despite worse baseline conditions).

Timeline:
\begin{itemize}
\item Week 1-2: Crisis shock; E drops to 5.6\%
\item Week 2-4: Implement redesigned portfolio
\item Month 2-3: Stabilization at 8-10\%
\item Month 4-6: Recovery to 12-14\% as new routine emerges
\item Month 9+: Convergence when restructuring is complete
\end{itemize}

Key lesson: Portfolio redesign requires understanding which interventions \textbf{backfire} in crisis (autonomy, recognition) and which \textbf{strengthen} (transparency, safety, support). One-size-fits-all fails.

% =============================================================================
% SECTION 5: CROSS-DOMAIN COHERENCE AND SPILLOVERS
% =============================================================================

\section{Cross-Domain Coherence: When Policies Conflict}
\label{sec:cross-domain}

\subsection{Policy Spillovers Between Domains}

So far we designed portfolios within single domains: Health, Pensions, Environment, Labor. But real policy operates across domains simultaneously.

A key finding: \textbf{Policies in one domain affect outcomes in other domains}.

Example: Pension reform (Domain 2) affects labor market behavior (Domain 4):

\begin{quote}
\textbf{Spillover:} If government implements flexible retirement pathway (WHEN choice architecture in Domain 2), this affects which workers stay in labor force. If combined with labor market policies encouraging longer working lives (Domain 4), the two policies \emph{reinforce}. But if they send contradictory messages (``retire early if you want'' vs. ``stay in workforce longer''), they \emph{interfere}.
\end{quote}

\subsection{Spillover Analysis Framework}

For any pair of domains with policies, analyze three questions:

\begin{enumerate}
\item \textbf{Congruence:} Do the policies send the same message to individuals?
\item \textbf{Capacity:} Do the policies compete for limited implementation resources?
\item \textbf{Behavioral Coherence:} Do they target the same population segments?
\end{enumerate}

Example: Health (Domain 1) + Environment (Domain 3):

\begin{center}
\begin{tabular}{@{}llll@{}}
\toprule
\textbf{Dimension} & \textbf{Health Policy} & \textbf{Environmental Policy} & \textbf{Assessment} \\
\midrule
Congruence & ``Reduce diabetes via lifestyle'' & ``Reduce energy via behavior change'' & Aligned (+0.4 γ) \\
Message & ``Take care of your health'' & ``Take care of the environment'' & Synergistic framing \\
Implementation & Health department leads & Environmental dept. leads & Parallel, not competing \\
Segment target & Autonomy-Seeking (health) & Social-Oriented (environment) & Different, no direct conflict \\
\midrule
\end{tabular}
\end{center}

Result: These portfolios have \textbf{positive spillover} (γ > 0). Combined messaging ``personal health and environmental stewardship go together'' amplifies each.

Contrast with Health + Pension:

\begin{center}
\begin{tabular}{@{}llll@{}}
\toprule
\textbf{Dimension} & \textbf{Health Policy} & \textbf{Pension Policy} & \textbf{Assessment} \\
\midrule
Congruence & ``Invest in health behaviors now'' & ``Work longer to fund pensions'' & Mixed (0.1 γ) \\
Message & ``Change lifestyle for health'' & ``Change retirement expectations'' & Competing frames \\
Capacity & Health dept implementation & Pension dept implementation & Parallel, but both need \\
& & & individual time (conflict) \\
Segment target & Health-motivated (D1) & Financial-motivated (D2) & Different motivations \\
\midrule
\end{tabular}
\end{center}

Result: These portfolios have \textbf{neutral/weak spillover} (γ ≈ 0). They don't interfere but also don't reinforce.

\subsection{Coherence Score for Multi-Domain Policy}

For a national government implementing policies across multiple domains simultaneously, compute a \textbf{Policy Coherence Score}:

$$\text{Coherence} = \frac{\sum_{i < j} \gamma_{ij}}{\text{# of domain pairs}}$$

Example calculation for 4 domains:

\begin{center}
\begin{tabular}{@{}llc@{}}
\toprule
\textbf{Domain Pair} & \textbf{Spillover Effect} & \textbf{γ} \\
\midrule
Health ↔ Pension & Conflicting messages & 0.1 \\
Health ↔ Environment & Aligned messages & 0.4 \\
Health ↔ Labor & Neutral (different targets) & 0.0 \\
Pension ↔ Environment & Weak interaction & 0.2 \\
Pension ↔ Labor & Strong: work-life interaction & 0.6 \\
Environment ↔ Labor & Weak: green jobs effect & 0.1 \\
\midrule
\textbf{Average Coherence Score} & & 0.23 \\
\midrule
\end{tabular}
\end{center}

Interpretation:
- Score > 0.3: Policies well-coordinated (reinforce each other)
- Score 0.1-0.3: Policies neutral (don't interfere much)
- Score < 0.1: Policies poorly coordinated (may interfere)

At 0.23, this policy mix is ``acceptable but could improve''. The weak point: Pension-Health spillover (γ=0.1) should be strengthened with better message alignment.

\subsection{Recommendations for Cross-Domain Coherence}

To improve the Policy Coherence Score:

\begin{enumerate}
\item \textbf{Strengthen Pension-Health alignment:} Frame pension flexibility as ``staying active and healthy in later years,'' not just ``financial necessity.'' This increases γ from 0.1 → 0.3.

\item \textbf{Amplify the Pension-Labor link:} Coordinate messaging so labor market policies (``skills training at age 55'') align with pension policies (``retirement flexibility at age 62''). This can increase γ from 0.6 → 0.7.

\item \textbf{Reduce message fatigue:} Population faces 4 domains of policy messages simultaneously. Batch related messages (health + environment) together; separate time-sensitive messages (pension reform) from ongoing ones (health awareness).

\item \textbf{Segment-specific integration:} Autonomy-Seeking segments respond to ``choice and control'' messaging (pension flexibility + health autonomy). Social-Oriented respond to ``community and fairness'' (pension intergenerational fairness + environmental stewardship). Message each segment separately.
\end{enumerate}

Revised Coherence Score after coordination: 0.23 → 0.35 (entering ``well-coordinated'' range).

% =============================================================================
% SECTION 6: IMPLEMENTATION GUIDANCE AND INTEGRATION
% =============================================================================

\section{Practical Implementation Framework}
\label{sec:implementation}

\subsection{Step-by-Step Process for Policy Designers}

For any new policy domain, follow this workflow:

\begin{enumerate}
\item \textbf{Diagnose Status Quo (Ch16 framework):}
   \begin{itemize}
   \item Measure 9C dimensions for target population
   \item Identify phase distribution (PreAware / Triggered / Action / Maintenance)
   \item Segment the population by behavioral type
   \item Document context dimensions (Ψ)
   \end{itemize}

\item \textbf{Design Single Interventions (Ch17 framework):}
   \begin{itemize}
   \item Let intervention vectors emerge from the 10C Status Quo profile
   \item Estimate baseline effect for each (literature + pilots)
   \item Document implementation constraints
   \end{itemize}

\item \textbf{Modulate by Phase and Segment (Ch18-19):}
   \begin{itemize}
   \item Build phase-affinity table: α for each intervention × phase combo
   \item Build segment-multiplier table: σ for each intervention × segment combo
   \item Identify which interventions have high phase-affinity for your population's current phase
   \item Identify which segments will be underserved or resistant
   \end{itemize}

\item \textbf{Optimize Static Portfolio (Ch20):}
   \begin{itemize}
   \item Select 3-5 high-affinity interventions
   \item Estimate complementarities (γ matrix)
   \item Apply coherence criteria (C1-C4 from Ch20)
   \item Compute combined effectiveness using Eq. 906
   \item Check implementation feasibility (budget, capacity, political)
   \end{itemize}

\item \textbf{Plan Dynamic Adjustment (Ch21):}
   \begin{itemize}
   \item Identify trigger conditions for phase progression (when does φ change?)
   \item Identify trigger conditions for context shift (when does Ψ change?)
   \item Identify trigger conditions for convergence (when to stop adapting?)
   \item Plan redesign actions for each trigger
   \item Set measurement timeline (quarterly review recommended)
   \end{itemize}

\item \textbf{Check Cross-Domain Coherence:}
   \begin{itemize}
   \item Identify other domains affected by your policy portfolio
   \item Compute spillover effects (γ between domains)
   \item Adjust messaging for alignment with other domains
   \item Target Coherence Score > 0.25 (acceptable) or > 0.35 (good)
   \end{itemize}

\end{enumerate}

\subsection{Key Metrics for Monitoring}

For ongoing implementation, track these metrics quarterly:

\begin{center}
\begin{tabular}{@{}llll@{}}
\toprule
\textbf{Metric} & \textbf{Formula} & \textbf{Target} & \textbf{Action if Miss} \\
\midrule
Awareness (A) & \% of population aware of policy & +20pp/year & Boost communication (AWARE) \\
Readiness (W) & \% ready to change behavior & +15pp/year & Check message clarity or barriers \\
Phase shift (φ) & \% moving to next phase & +10-15pp/year & Assess if interventions matched phase \\
Effectiveness (E) & Measured behavior change & Project target & Redesign if <80\% of projection \\
Convergence & ||Θ\_t - Θ_{t-1}|| & <0.05/quarter & If persistent change, learn and update \\
Coherence & Cross-domain spillovers & >0.25 & Adjust messaging across domains \\
\midrule
\end{tabular}
\end{center}

\subsection{Common Pitfalls to Avoid}

\begin{enumerate}
\item \textbf{Ignoring phase distribution:} Applying WHAT$_F$ (financial incentives) to PreAware population. Use AWARE-READY first.

\item \textbf{One-size-fits-all design:} Treating all segments identically. Autonomy-Seeking hate mandates; Social-Oriented seek fairness narratives. Customize.

\item \textbf{Skipping complementarity analysis:} Adding interventions without checking γ matrix. Some combinations backfire.

\item \textbf{Static thinking:} Designing portfolio for today, not accounting for phase progression or context shifts. Plan redesigns ahead.

\item \textbf{Implementation overload:} Trying to implement 8+ interventions simultaneously. Start with 4-5, add strategically as capacity grows.

\item \textbf{Neglecting spillovers:} Designing domain-by-domain without checking coherence with other policies. Can create contradiction and backlash.

\item \textbf{Measurement gaps:} Not measuring parameters as you learn. Miss opportunity to update effectiveness estimates (Bayesian learning from Ch21).

\end{enumerate}

% =============================================================================
% READING PATHS AND NAVIGATION
% =============================================================================

\section*{Reading Paths for This Chapter}

\begin{tcolorbox}[colback=yellow!5!white, colframe=yellow!75!black, title=Where to Go Next]

\textbf{If you want...} then:
\begin{itemize}[nosep]
\item To apply framework to your specific domain → Pick Section 1-4 matching your domain
\item Copy-paste templates for your domain → Appendix UNMAPPED_LLL (templates + matrices)
\item To understand domain interactions → Section 5 (Cross-Domain Coherence)
\item To analyze spillovers & optimize coherence → Appendix UNMAPPED_MMM (spillover framework)
\item To implement a policy portfolio → Section 6 (Practical Implementation)
\item To scale across multiple organizational levels → Chapter 23 (Multi-Level Implementation)
\item Formal mathematics behind portfolio dynamics → Appendix UNMAPPED_KKK
\item Parameter sourcing discipline for your domain → Appendix UNMAPPED_JJJ
\item Case studies of policy successes/failures → Case studies in Appendix UNMAPPED_MMM
\end{itemize}

\end{tcolorbox}

% =============================================================================
% SUMMARY SECTION
% =============================================================================
\section{Summary}
\label{sec:ch22-summary}

\begin{tcolorbox}[
    colback=blue!5!white,
    colframe=blue!75!black,
    title={\textbf{Chapter 22 Summary: Policy Applications}}
]

\textbf{Core Insight:} Policy portfolio design follows the same principles as individual-level design, but with scale-up of segment heterogeneity, implementation dynamics, and cross-domain spillovers.

\textbf{Key Results:}
\begin{enumerate}[nosep]
    \item \textbf{Health Portfolio:} PreAware populations need AWARE emphasis; $K^*_{\text{health}} \approx 4$ (L1 scope)
    \item \textbf{Pension Reform:} Context-shifting requires behavioral + institutional; $K^*_{\text{pension}} \approx 3$ (L0 scope)
    \item \textbf{Environmental Policy:} Crisis contexts trigger $K^*(t) \uparrow$; redesign with WHAT$_F$ emphasis
    \item \textbf{Labor Markets:} Disruption collapses E by 64\%; $K^*$ increases temporarily for recovery
    \item \textbf{Cross-Domain:} Policy Coherence Score $> 0.25$; C2 constraint: $\gamma_{ij} \geq -0.1$
    \item \textbf{K* Policy Pattern (EIT-EMP-2):} L0/L1 scope $\to$ $K^* \approx 3-6$ (fewer, deeper interventions)
\end{enumerate}

\textbf{Practical Framework (mit EMERGE):}
\begin{center}
\small
Ch11 (EMERGE, $K^*$) $\rightarrow$ Ch16 ($S_0$) $\rightarrow$ Ch17 ($\vec{I}$) $\rightarrow$ Ch18-19 ($\alpha, \sigma$) $\rightarrow$ Ch20 (C1-C7) $\rightarrow$ Ch21 ($K^*(t)$) $\rightarrow$ \textbf{Ch22 (Policy)}
\end{center}

\textbf{Policy-Level Scope:} L0 (Systemic) und L1 (Strategic) mit $f_L \in [0.6, 0.8]$

\end{tcolorbox}

\subsection{Formal Details}
\label{sec:ch22-formal}

The core equations underlying this chapter:

\textbf{Portfolio Effectiveness (from Ch20):}
\begin{equation}
\label{eq:ch22-portfolio}
E(P) = \sum_{i} E_i \cdot \alpha_i \cdot \sigma_s + \sum_{i<j} \gamma_{ij} \sqrt{E_i \cdot E_j}
\end{equation}
where $E_i$ is single intervention effect, $\alpha_i$ is phase-affinity, $\sigma_s$ is segment multiplier, and $\gamma_{ij}$ is complementarity.

\textbf{Policy Coherence Score:}
\begin{equation}
\label{eq:ch22-coherence}
\text{Coherence} = \frac{\sum_{i < j} \gamma_{ij}}{\binom{n}{2}}
\end{equation}
where $n$ is number of policy domains.

\textbf{Crisis Adjustment Factor:}
\begin{equation}
\label{eq:ch22-crisis}
E_{\text{crisis}} = E_{\text{base}} \cdot \kappa(\Psi) \quad \text{where } \kappa(\Psi) \in [0.4, 1.0]
\end{equation}

\textbf{Policy-Level K* (EIT-EMP-2, EXC-3 Hybrid):}
\begin{equation}
\label{eq:ch22-k-policy}
K^*_{\text{policy}} = K_{\text{base}} \cdot f_B + \delta_C + \delta_S + \delta_\gamma + \delta_L \quad \text{where } \delta_L \in [-40\%, -20\%] \cdot K_{\text{base}} \text{ for L0/L1}
\end{equation}

\subsection{What Comes Next}
\label{sec:ch22-next}

This chapter demonstrated policy application in single organizations or regions. The next chapters extend the framework:

\begin{enumerate}[nosep]
    \item \textbf{Chapter 23:} Multi-level implementation---scaling from individual to organizational to societal levels
    \item \textbf{Chapter 24:} Emergent life journeys---intergenerational dynamics and lifespan-level patterns
\end{enumerate}

%======================================================================
% END OF CHAPTER 22
%======================================================================

\newpage
